%=========== ΛΥΣΗ - 40ο ΘΕΜΑ Α ===========
\providecommand{\theoriaref}[3]{%
\begin{center}
\fcolorbox{maincolor!70!black}{maincolor!8}{%
\begin{minipage}{.88\linewidth}
\textcolor{maincolor!80!black}{\kerkissans{\textbf{#1~\ref{#2}}}}\\[1mm]
#3
\end{minipage}}
\end{center}}
\providecommand{\orismosref}[1]{%
\theoriaref{Ορισμός}{#1}{Η απάντηση δίνεται από τον αντίστοιχο ορισμό.}}
\providecommand{\apodeixiref}[1]{%
\theoriaref{Θεώρημα}{#1}{Η ζητούμενη απόδειξη δίνεται στην αντίστοιχη απόδειξη.}}

\begin{thema}{Α}
\begin{erwthma}
\item \apodeixiref{thewrhma:paragwgos_synefaptomenhs}

\item \orismosref{orismos:grafikh_parastash}

\item \orismosref{orismos:rythmos_metavolhs}

\item Οι χαρακτηρισμοί των προτάσεων είναι:
\begin{alist}
\item \textbf{Σωστό}. Η ποσότητα $f(a)$ είναι σταθερή ως προς τη μεταβλητή ολοκλήρωσης $x$, άρα
\[ \dint_a^\beta f(a)\d x=f(a)(\beta-a). \]
\item \textbf{Σωστό}. Αν μια συνάρτηση ήταν παραγωγίσιμη στο $x_0$, τότε θα ήταν και συνεχής στο $x_0$. Άρα, αν δεν είναι συνεχής, δεν είναι παραγωγίσιμη.
\item \textbf{Σωστό}. Το εμβαδόν χωρίου είναι γεωμετρικό μέγεθος και είναι πάντοτε μη αρνητικός αριθμός.
\item \textbf{Λάθος}. Για να είναι μια συνάρτηση παράγουσα κάποιας άλλης, πρέπει να είναι παραγωγίσιμη. Δεν είναι κάθε συνεχής συνάρτηση παραγωγίσιμη.
\item \textbf{Λάθος}. Το θεώρημα \eng{Bolzano} απαιτεί συνέχεια στο $[a,\beta]$ και ετερόσημες τιμές στα άκρα, όχι γνήσια μονοτονία.
\end{alist}

\item Σωστή απάντηση είναι η \textbf{β}. Από το σχήμα έχουμε $f'(x)>0$ για $x<3$ και $f'(x)<0$ για $x>3$, οπότε η $f$ παρουσιάζει τοπικό μέγιστο στη θέση $x=3$.
\end{erwthma}
\end{thema}
%=========== ΤΕΛΟΣ ΛΥΣΗΣ - 40ο ΘΕΜΑ Α ===========
